\documentclass[10pt,a4paper]{amsart}
\usepackage[margin=19mm]{geometry}
\usepackage{amsmath,amssymb,mathtools}
\usepackage[scr=rsfs,cal=euler]{mathalfa}
\usepackage{xfrac}
\usepackage[unicode,hidelinks,bookmarks=false]{hyperref}
\newcommand{\ie}{i.e.\ }
\newcommand{\cf}{cf.\ }
\newcommand{\resp}{respectively\ }
\newcommand{\Subsection}{\subsection}
\providecommand{\nobreakdash}{\nobreak\hbox{-}}
\setlength{\emergencystretch}{2em}
\multlinegap0pt
\usepackage{bm}
\usepackage{bbm}
\usepackage{dsfont}
\usepackage{xspace}
\usepackage{cancel}
\usepackage{mathtools}
\usepackage{amsthm}
\makeatletter
\@ifundefined{theorem}{\newtheorem{theorem}{Theorem}}{}
\makeatother
\title[Non-Binary Quasi-Cyclic LDPC Codes with Entanglement Assista]{Non-Binary Quasi-Cyclic LDPC Codes with Entanglement Assistance}
\author{Pavan Kumar, Shayan Srinivasa Garani}
\date{Source: arXiv:2608.16761v1 (CC BY 4.0)}
\begin{document}
\maketitle

\noindent\emph{Source and licence.} Non-Binary Quasi-Cyclic LDPC Codes with Entanglement Assistance. Version arXiv:2608.16761v1 (CC BY 4.0), distributed under the Creative Commons Attribution 4.0 International licence, as recorded in the arXiv licence metadata for this exact version and shown on the abstract page https://arxiv.org/abs/2608.16761v1. Full rights notice: licenses/arxiv-2608.16761-CC-BY-4.0.txt.\par\smallskip

\noindent\textit{Editorial scope.} Counted unit: the Theorem whose source label is \texttt{theorem1} in the article (source0.tex lines 354--356, source label \texttt{theorem1}), delivered together with the article's own complete original proof of it (source0.tex lines 357--434, one \texttt{proof} environment of 78 lines). No mathematics is written, repaired or re-derived: statement and proof are the article's own text line for line. Required context retained uncounted: sumofrowzero (lines 113--115). Every definition, display and label of those lines is retained unchanged. Omitted from this package and not redistributed: 1--112, 116--353, 435--642. Nothing inside a retained excerpt is omitted or altered: every retained body is delivered exactly as the article's lines are written, so no omission receipt of any kind accompanies this package. Citations: the retained excerpts carry no citation command at all, so no bibliography entry of the article is transcribed and no reference is redistributed. Reference adaptation: none was needed and none was made; every reference inside the delivered unit resolves inside it, so no unresolved reference or citation is produced. Printed slips kept literal: no printed word, symbol, hypothesis or number of the counted statement or of its proof was altered. Layout and class: the article's own class is \texttt{IEEEtran} with packages including footnote, inputenc, babel, fontenc, amsmath, amssymb, amsthm, mathtools, braket, halloweenmath, cite, graphicx, xcolor, tcolorbox; the wrapper is the lane's pinned \texttt{amsart} editorial wrapper at 10pt on A4, which restates the theorem-like environments the retained text needs and adds the article's own macro definitions that the retained text needs (none), with extra wrapper packages (bm, bbm, dsfont, xspace, cancel, mathtools). The delivered numbering is the unit's own. Assets: no retained span contains a figure environment, a graphics-inclusion command, a PDF-page inclusion, a table or a TikZ picture; no image file is copied into this package, the package declares no asset dependency and no image file is redistributed with it. Licence: version arXiv:2608.16761v1 of this article is distributed under the Creative Commons Attribution 4.0 International licence, as recorded in the arXiv licence metadata for this exact version and shown on the abstract page https://arxiv.org/abs/2608.16761v1. A full rights notice is delivered at licenses/arxiv-2608.16761-CC-BY-4.0.txt. Characters: every retained excerpt is pure ASCII, so the delivered engine, XeLaTeX with its default Latin Modern text font, reports no missing character.
\begin{equation}\label{sumofrowzero}
  \sum_{k=0}^{n-1} a_k^{-1} R_k = \mathbf{0},
\end{equation}

\begin{theorem}\label{theorem1}
	Let $H$ be a submatrix of $\mathbf{H}$ consisting distinct $k$ $(1\leq k\leq r)$ block-rows. Then, $\mathrm{gfrank}(H)=r+(k-1)(r-1)$ over $\mathbb{F}_{q}$.
\end{theorem}

\begin{proof} Suppose  $k = 1.$ In this case, $\mathrm{gfrank}_{q}(H) = r$, as the matrix is already in row echelon form.

Next, consider the case when \( k = 2 \). Subtracting the first block-row from the second and applying \eqref{sumofrowzero} results in the last row of the second block-row becoming zero, i.e., any two block-rows  of $H$ are linearly dependent.

Now, consider the $k$ block-rows of $H$, and remove the last row from each block-row except the first. Let
\begin{equation} \label{linearcombi}
			\sum_{i_{1}=0}^{r-1}\alpha_{i_{1}}^{(0)}R^{(0)}_{i_{1}}+\sum_{i_{2}=0}^{r-2}\alpha_{i_{2}}^{(1)}R^{(1)}_{i_{2}}+\cdots+\sum_{i_{k}=0}^{r-2}\alpha_{i_{k}}^{(k-1)}R^{(k-1)}_{i_{k}}=0,
		\end{equation}
        where coefficients are from the field $\mathbb{F}_{q}$. 
        
        Define the inner product $\langle,\rangle:\mathbb{F}_{q}^{r^{2}}\times \mathbb{F}_{q}^{r^{2}}\to \mathbb{F}_{q}$, as follows:		\begin{equation}\label{euclideanF}		\langle(v_{1},\ldots,v_{r^{2}}),(u_{1},\ldots,u_{r^{2}})\rangle=\sum_{i=1}^{r^{2}}v_{i}u_{i}.
	\end{equation}
Taking the inner product, as defined by \eqref{euclideanF}, of \eqref{linearcombi} with each of the remaining rows from all the block-rows  $R^{(0)}, R^{(1)}, \dots, R^{(k-1)}$, and applying Lemma 1, results in the following system of linear equations:	
	\begin{align*}
	& r\alpha_{j_{1}}^{(0)}(a_{j_{1}})^{2}+\sum\limits_{i_{2}=0}^{r-2}\alpha_{i_{2}}^{(1)}a_{i_{2}}a_{j_{1}}+\cdots+\sum\limits_{i_{k}=0}^{r-2}\alpha_{i_{k}}^{(k-1)}a_{i_{k}}a_{j_{1}}=0,\\
    &\text{ for all } 0\leq j_{1}\leq r-1;\\
	&\sum\limits_{i_{1}=0}^{r-1}\alpha_{i_{1}}^{(0)}a_{i_{1}}a_{j_{2}}+r\alpha_{j_{2}}^{(1)}(a_{j_{2}})^{2}+\cdots+\sum\limits_{i_{k}=0}^{r-2}\alpha_{i_{k}}^{(k-1)}a_{i_{k}}a_{j_{2}}=0,\\
    &\text{ for all } 0\leq j_{2}\leq r-2;\\
	&\hspace{44.5mm}\vdotswithin{=}\\	
	&\sum\limits_{i_{1}=0}^{r-1}\alpha_{i_{1}}^{(0)}a_{i_{1}}a_{j_{k}}+\sum\limits_{i_{2}=0}^{r-2}\alpha_{i_{2}}^{(1)}a_{i_{2}}a_{j_{k}}+\cdots+r\alpha_{j_{k}}^{(k-1)}(a_{j_{k}})^{2}=0,\\
    &\text{ for all } 0\leq j_{k}\leq r-2.
	\end{align*}	
    Let  $X^{(i)}=[\alpha^{(i)}_{0},\alpha^{(i)}_{1},\ldots,\alpha^{(i)}_{(r-2)}]^{T}$, for all {\small$1\leq i\leq k-1$} and $X^{(0)}=[\alpha^{(0)}_{0},\alpha^{(0)}_{1},\ldots,\alpha^{(0)}_{(r-1)}]^{T}$, and let
    \begin{align*}
        A&=[a_{ij}], \text{where } a_{ij}=ra_{i}^{2}\delta_{ij}, 0\leq i,j\leq r-1;\\
        B&=[b_{ij}], \text{where } b_{ij}=ra_{i}^{2}\delta_{ij}, 0\leq i,j\leq r-2;\\
        C&=[c_{ij}], \text{where }  c_{ij}=a_{i}a_{j},  0\leq i\leq r-1, 0\leq j\leq r-2;\\
        D&=[d_{ij}], \text{where }  d_{ij}=a_{i}a_{j}, 0\leq i\leq r-2,0\leq j\leq r-1;\\
        E&=[e_{ij}], \text{where } e_{ij}=a_{i}a_{j}, 0\leq i,j\leq r-2.
    \end{align*}
   The above system of linear equations can be written as follows:
   {\small \begin{equation}
       \begin{bmatrix}
           A&C&C&\cdots&C\\
           D&B&E&\cdots&E\\
           D&E&B&\cdots&E\\
           \vdots&\vdots&\vdots&\ddots&\vdots\\
           D&E&E&\cdots&B
       \end{bmatrix}\begin{bmatrix}
           X^{(0)}\\
           X^{(1)}\\
           X^{(2)}\\
           \vdots\\
            X^{(k-1)}
       \end{bmatrix}=\begin{bmatrix}
           0\\
           0\\
           0\\
           \vdots\\
            0
       \end{bmatrix}.
   \end{equation}}
	It is straightforward to check that 
   {\small \begin{equation}\label{deteq}
        \begin{vmatrix}
           A&C&C&\cdots&C\\
           D&B&E&\cdots&E\\
           D&E&B&\cdots&E\\
           \vdots&\vdots&\vdots&\ddots&\vdots\\
           D&E&E&\cdots&B 
        \end{vmatrix}= a_{r-1}\prod_{i=0}^{r-2}a_{i}^{k}\begin{vmatrix}
           A'&C'&C'&\cdots&C'\\
           D'&B'&E'&\cdots&E'\\
           D'&E'&B'&\cdots&E'\\
           \vdots&\vdots&\vdots&\ddots&\vdots\\
           D'&E'&E'&\cdots&B' 
        \end{vmatrix},
    \end{equation}}
    where
    \vspace{-0.25cm}
    \begin{align*}
         A'&=[a'_{ij}], \text{where } a'_{ij}=ra_{i}\delta_{ij}, 0\leq i,j\leq r-1;\\
          B'&=[b'_{ij}], \text{where } b'_{ij}=ra_{i}\delta_{ij},0\leq i,j\leq r-2;\\
           C'&=[c'_{ij}], \text{where } c'_{ij}=a_{j},0\leq i\leq r-1,0\leq j\leq r-2;\\
             D'&=[d'_{ij}], \text{where } d'_{ij}=a_{j}, 0\leq i\leq r-2,0\leq j\leq r-1;\\
                E'&=[e'_{ij}], \text{where } e'_{ij}=a_{j},  0\leq i,j\leq r-2.
    \end{align*}
By applying elementary row operations, we can verify that the determinant on the right-hand side of \eqref{deteq} is nonzero. This implies that the system of linear equations has only the zero vector as a solution, completing the proof.\end{proof}

\end{document}
